% ===========================================================================
\chapter{A geometry toolkit}\label{ch:geometry}

With everything on a flat plane in metres, a map needs only a handful of
geometric tools, and they are all in this chapter:

\begin{itemize}[nosep]
  \item the length of a road, and the point a given distance along it;
  \item the point of a road closest to a GPS fix, and how far away it is;
  \item the area of a lake, and whether a point is inside a wood;
  \item which way a road points, and how sharply it turns.
\end{itemize}

\section{Lines}

On a map, a road or a river is a \emph{polyline}: a list of points joined
by straight \emph{segments}. Its length is the sum of the lengths of its
segments.

\begin{newfile}{geometry.nex}
-- geometry.nex: the flat geometry a map needs. Lines, segments and
-- polygons, and the directions of travel along them.

intern Geo

-- The total length of a line through `points`, in order.
function line_length(points: Array[Point]): Real do
  from
    let i: Integer := 1
  until
    i >= points.length
  do
    result := result + points.get(i - 1).distance_to(points.get(i))
    i := i + 1
  end
ensure
  non_negative: result >= 0.0
end
\end{newfile}

\code{result} starts at \code{0.0}, the zero value of \code{Real}, so the
loop can add to it straight away.

To find the point \code{d} metres along a line, walk along it segment by
segment, subtracting each segment's length from \code{d}, until \code{d}
is less than the length of the segment you are on. The point is then that
far into the segment: a fraction \code{d / length} of the way from its
start \code{a} to its end \code{b}, which is \code{a + (b - a) * (d / length)}.

\begin{append}{geometry.nex}
-- The point `d` metres along the line through `pts` (its last point if d
-- is past the end).
function point_along(pts: Array[Point], d: Real): Point
require
  not_empty: not pts.is_empty
do
  result := pts.get(pts.length - 1)
  let left: Real := d
  let found: Boolean := false
  from
    let i: Integer := 1
  until
    found or i >= pts.length
  do
    let a: Point := pts.get(i - 1)
    let b: Point := pts.get(i)
    let seg: Real := a.distance_to(b)
    if left <= seg and seg > 0.0 then
      result := a + (b - a) * (left / seg)
      found := true
    else
      left := left - seg
    end
    i := i + 1
  end
end

-- The average of some points: their centre of mass.
function average(pts: Array[Point]): Point
require
  not_empty: not pts.is_empty
do
  let sum: Point := create Point.make(0.0, 0.0)
  across pts as p do
    sum := sum + p
  end
  result := sum * (1.0 / pts.length.to_real)
end
\end{append}

The average of the corners of an area is a good enough place to write the
area's name on a map; we'll use it for exactly that.

\section{Segments, and the dot product}

When a GPS says you are at some point $p$, a navigator wants to know which
point of the road you are on: the point of the segment from $a$ to $b$
closest to $p$.

\begin{center}
\begin{tikzpicture}[scale=1.1]
  \coordinate (A) at (0,0);
  \coordinate (B) at (6,1);
  \coordinate (P) at (2.6,2.2);
  \coordinate (C) at ($(A)!(P)!(B)$);
  \draw[accent, very thick] (A) -- (B);
  \draw[color=cm, dashed] (P) -- (C);
  \draw[color=cm, -{Stealth}] (A) -- (P);
  \fill (A) circle (0.06) node[below left] {$a$};
  \fill (B) circle (0.06) node[below right] {$b$};
  \fill[st] (P) circle (0.07) node[above] {$p$};
  \fill[accent] (C) circle (0.07) node[below] {closest point};
  \draw[color=cm] ($(C)!0.25cm!(B)$) -- ($($(C)!0.25cm!(B)$)!0.25cm!90:(B)$) -- ($(C)!0.25cm!(P)$);
\end{tikzpicture}
\end{center}

Every point of the segment is $a + t\,(b - a)$ for some $t$ from 0 (at $a$)
to 1 (at $b$). The closest one is where the dashed line from $p$ meets the
segment at a right angle, and the dot product finds it. For two arrows
$u$ and $v$,
\[
  u \cdot v = u_x v_x + u_y v_y = |u|\,|v|\cos\theta,
\]
where $\theta$ is the angle between them. So $u \cdot v / |v|$ is
$|u|\cos\theta$: how far $u$ reaches in the direction of $v$, the length of
its \emph{shadow} on $v$. With $u = p - a$ and $v = b - a$, the shadow is
how far along the segment the closest point is, and dividing by the
segment's length once more turns it into the fraction we want:
\[
  t = \frac{(p - a) \cdot (b - a)}{|b - a|^2}.
\]
If $t$ comes out below 0 or above 1, the right angle falls off the end of
the segment, and the nearest point is the nearer end. \code{clamp(0.0, 1.0)}
pulls $t$ back into range, which picks exactly that end.

\begin{append}{geometry.nex}
-- The point of the segment a-b closest to p.
--
-- Every point of the segment is a + t (b - a) for some t between 0 and 1.
-- The closest one to p is where the line from p meets the segment at a
-- right angle, which is at t = (p - a)·(b - a) / |b - a|^2. If that t is
-- below 0 or above 1, the right angle is off the end of the segment, and
-- the closest point is the nearer end.
function closest_point_on_segment(p, a, b: Point): Point do
  let ab: Point := b - a
  let len2: Real := ab.dot(ab)
  if len2 = 0.0 then
    result := a
  else
    let t: Real := ((p - a).dot(ab) / len2).clamp(0.0, 1.0)
    result := a + ab * t
  end
ensure
  no_further_than_ends: p.distance_to(result) <= p.distance_to(a) + 0.000001
    and p.distance_to(result) <= p.distance_to(b) + 0.000001
end

function distance_to_segment(p, a, b: Point): Real do
  result := p.distance_to(closest_point_on_segment(p, a, b))
end
\end{append}

The postcondition is a sanity check we can state without knowing anything
about dot products: the closest point can be no further away than either
end. The tiny allowance, a millionth of a metre, is for rounding.

\section{Polygons}

An area, such as a lake or a wood, is a \emph{polygon}: its outline is a
list of corners, and the last corner joins back to the first.

\subsection{Area: the shoelace formula}

The area of a polygon with corners $(x_0, y_0), \dots, (x_{n-1}, y_{n-1})$ is
\[
  A = \frac{1}{2}\,\Big|\sum_{i=0}^{n-1} (x_i\, y_{i+1} - x_{i+1}\, y_i)\Big|,
\]
where the corner after the last is the first again. It is called the
\emph{shoelace formula} because of the criss-cross pattern of the products
when the corners are written in a column.

Why does it work? Each term $x_i y_{i+1} - x_{i+1} y_i$ is twice the area
of the triangle from the origin to one edge of the outline, positive if
the edge goes anticlockwise around the origin and negative if it goes
clockwise. Going all the way round the outline, the triangles outside the
polygon are counted once each way and cancel, and what is left is the
polygon itself. The absolute value at the end makes the answer positive
whichever way round the corners are listed.

\begin{append}{geometry.nex}
-- The area inside a polygon, in square metres, by the shoelace formula.
-- `ring` lists the corners in order around the outline; the last corner
-- joins back to the first.
function polygon_area(ring: Array[Point]): Real
require
  is_polygon: ring.length >= 3
do
  let twice: Real := 0.0
  from
    let i: Integer := 0
  until
    i = ring.length
  do
    let p: Point := ring.get(i)
    let q: Point := ring.get((i + 1) % ring.length)
    twice := twice + (p.x * q.y - q.x * p.y)
    i := i + 1
  end
  result := twice.abs / 2.0
ensure
  non_negative: result >= 0.0
end
\end{append}

\texttt{(i + 1) \% ring.length} is the index of the next corner, which is
0 again after the last one.

\subsection{Inside or outside: casting a ray}

Is a point inside a polygon? Walk from the point in a straight line,
due east, all the way to the edge of the world, and count how many times
you cross the outline. Every crossing takes you from inside to outside or
back, and you end up outside. So if you crossed an odd number of times, you
started inside; an even number, outside.

\begin{center}
\begin{tikzpicture}[scale=0.9]
  \fill[titlebg, draw=accent] (0,0) -- (3,-0.5) -- (5,1) -- (3.5,1.6) -- (4.6,3)
    -- (1.5,3.2) -- (-0.4,1.8) -- cycle;
  \fill[st] (1.2,1.2) circle (0.07) node[below] {inside};
  \draw[st, -{Stealth}] (1.2,1.2) -- (7,1.2);
  \fill[st] (4.4,1.2) circle (0.05);
  \node[font=\footnotesize, st] at (4.6,0.85) {1 crossing};
  \fill[red!70!black] (2.2,2.4) circle (0.07) node[above left] {inside};
  \draw[red!70!black, -{Stealth}] (2.2,2.4) -- (7,2.4);
  \node[font=\footnotesize, text=red!70!black] at (6.1,2.65) {3 crossings};
  \fill[color=cm] (5.4,0.3) circle (0.07) node[below] {outside};
  \draw[color=cm, -{Stealth}] (5.4,0.3) -- (7,0.3);
  \node[font=\footnotesize, text=cm] at (6.3,0.05) {0};
\end{tikzpicture}
\end{center}

For each edge from $a$ to $b$, the ray crosses it if one end is above the
point and the other is not, and the edge passes the point's height to the
right of the point. Where along the edge does it pass that height? A
fraction $(p_y - a_y)/(b_y - a_y)$ of the way from $a$ to $b$, so at
$x = a_x + \frac{p_y - a_y}{b_y - a_y}(b_x - a_x)$.

\begin{append}{geometry.nex}
-- Whether p is inside the polygon, by casting a ray from p due east and
-- counting how many edges of the outline it crosses: an odd number means
-- inside, an even number outside.
function polygon_contains(ring: Array[Point], p: Point): Boolean
require
  is_polygon: ring.length >= 3
do
  let inside: Boolean := false
  from
    let i: Integer := 0
  until
    i = ring.length
  do
    let a: Point := ring.get(i)
    let b: Point := ring.get((i + 1) % ring.length)
    -- Does edge a-b cross the horizontal line through p, to p's right?
    if (a.y > p.y) /= (b.y > p.y) then
      let crossing_x: Real := a.x + (p.y - a.y) / (b.y - a.y) * (b.x - a.x)
      if crossing_x > p.x then
        inside := not inside
      end
    end
    i := i + 1
  end
  result := inside
end
\end{append}

The test \code{(a.y > p.y) /= (b.y > p.y)} compares two \code{Boolean}s:
it is true when exactly one end of the edge is above the point. It also
guarantees that \code{b.y - a.y} isn't zero, so the division is safe.

\section{Directions}

Navigators, compasses and sailors give directions as \emph{bearings}:
angles in degrees measured clockwise from north, so that north is
$0^\circ$, east $90^\circ$, south $180^\circ$ and west $270^\circ$.

Mathematicians measure angles the other way: anticlockwise from east (the
$x$ axis). The function $\operatorname{atan2}(y, x)$ gives that angle for
the arrow $(x, y)$, correctly in all four quarters of the circle. Swapping
its arguments swaps the roles of the two axes, which turns the
mathematicians' angle into the compass's: $\operatorname{atan2}(x, y)$
measures from north (the $y$ axis) clockwise towards east. In Nex,
\code{y.atan2(x)} is $\operatorname{atan2}(y, x)$, so the bearing of the
arrow $d$ is \code{d.x.atan2(d.y)}. It comes out between $-180^\circ$ and
$180^\circ$, and we add $360^\circ$ to the negative ones.

\begin{append}{geometry.nex}
-- The compass bearing from a to b, in degrees: 0 is north, 90 east, 180
-- south and 270 west, measured clockwise as on a compass.
function bearing(a, b: Point): Real do
  let d: Point := b - a
  result := d.x.atan2(d.y).to_degrees
  if result < 0.0 then
    result := result + 360.0
  end
ensure
  in_range: result >= 0.0 and result < 360.0
end

-- An arrow of length 1 pointing along a compass bearing.
function direction(heading: Real): Point do
  let r: Real := heading.to_radians
  result := create Point.make(r.sin, r.cos)
end
\end{append}

\code{direction} goes the other way, from a bearing back to an arrow. North
is $(0, 1)$ and east is $(1, 0)$, so the arrow is
$(\sin\theta, \cos\theta)$, again with sine and cosine in the opposite
places to the usual $(\cos\theta, \sin\theta)$.

When a traveller heading on bearing \code{before} turns onto a road with
bearing \code{after}, how far do they turn? The difference
\code{after - before}, but measured the short way round: heading
$350^\circ$ and turning to $10^\circ$ is a turn of $20^\circ$ to the
right, not $340^\circ$ to the left. So a difference over $180^\circ$ has
$360^\circ$ taken off, and one of $-180^\circ$ or less has $360^\circ$
added. Positive turns are to the right (clockwise), negative ones to the
left.

\begin{append}{geometry.nex}
-- How far a traveller turns when changing from heading `before` to
-- heading `after`: positive to the right, negative to the left, and never
-- more than half a turn either way.
function turn_angle(before, after: Real): Real do
  result := after - before
  if result > 180.0 then
    result := result - 360.0
  elseif result <= -180.0 then
    result := result + 360.0
  end
ensure
  in_range: result > -180.0 and result <= 180.0
end

-- The nearest of the eight compass points to a bearing: "north",
-- "north-east", ..., "north-west".
function compass_name(degrees: Real): String
require
  in_range: degrees >= 0.0 and degrees < 360.0
do
  let names: Array[String] := ["north", "north-east", "east", "south-east",
    "south", "south-west", "west", "north-west"]
  result := names.get(((degrees / 45.0).round) % 8)
end
\end{append}

The eight compass points are $45^\circ$ apart, so \code{degrees / 45.0},
rounded, is the number of the nearest one. A bearing of $350^\circ$ rounds
to 8, one past north-west, and \texttt{\% 8} turns that back into 0, north.

\begin{tryit}{try_geometry.nex}
intern Geometry

function pt(x, y: Real): Point do
  result := create Point.make(x, y)
end

let road: Array[Point] := [pt(0.0, 0.0), pt(300.0, 0.0), pt(300.0, 400.0)]
console.print_line("the road is " + fixed(line_length(road), 0) + " m long")
console.print_line("200 m along it: " + point_along(road, 200.0).to_string())
console.print_line("500 m along it: " + point_along(road, 500.0).to_string())
let fix: Point := pt(120.0, 9.0)
console.print_line("closest point to " + fix.to_string() + ": "
  + closest_point_on_segment(fix, road.get(0), road.get(1)).to_string()
  + ", " + fixed(distance_to_segment(fix, road.get(0), road.get(1)), 1) + " m away")

let field: Array[Point] := [pt(0.0, 0.0), pt(100.0, 0.0), pt(100.0, 50.0), pt(0.0, 50.0)]
console.print_line("the field is " + fixed(polygon_area(field), 0) + " square metres")
console.print_line("(50, 25) inside? " + polygon_contains(field, pt(50.0, 25.0)).to_string)
console.print_line("(150, 25) inside? " + polygon_contains(field, pt(150.0, 25.0)).to_string)

let heading: Real := bearing(road.get(0), road.get(1))
let turned: Real := bearing(road.get(1), road.get(2))
console.print_line("first heading " + fixed(heading, 0) + " (" + compass_name(heading) + "), then "
  + fixed(turned, 0) + " (" + compass_name(turned) + ")")
console.print_line("the turn: " + fixed(turn_angle(heading, turned), 0) + " degrees")
\end{tryit}
\tryoutput{try_geometry}

The road heads east, then turns north: a turn of $-90^\circ$, a left turn.
\code{to_string} without parentheses is how Nex writes a call with no
arguments on a built-in value; \code{to_string()} works too.
